% !TEX program = xelatex
% !TEX encoding = UTF-8 Unicode
%
% <Course name> --- lecture notes, English edition.  One continuous document,
% ordered by topic.  Everything shared with the Korean edition is in
% coursenotes.sty and figures.tex.
% Build: ./build.sh   (latexmk -xelatex, three passes because of the ToC)
%
% Rules for adding a lecture (the same list heads notes-ko.tex):
%   1. Put new material in the part whose topic it belongs to.  Open a new part
%      with \partdivider only when a genuinely new topic starts.
%   2. Section numbers run straight through all parts; never restart them, and
%      never renumber a section that has been published (the site viewer keys
%      on them).
%   3. Put \lecturebreak[YYYY-MM-DD]{Lecture N} where that lecture's new
%      material begins, then a \sourcenote saying what it rests on, and add one
%      row to the lecture log below.
%   4. Keep sections, equations, figures, theorem-like environments and every
%      \label identical in order and name in the two editions.
%   5. Update \lecdate and the "arc of the course" list.
%
\documentclass[11pt]{article}
\usepackage{coursenotes}

%% words inside figures (the Korean edition defines the same macros)
\newcommand{\lblPosition}{position $x$}
\newcommand{\lblConcentration}{$c(x,t)$}
\newcommand{\lblSpreads}{width grows\\like $\sqrt{Dt}$}
%% figures.tex --- every diagram of the notes, shared by both editions.
%%
%% LANGUAGE-NEUTRAL ON PURPOSE.  No word is typed inside a figure: every label
%% goes through a \lbl... macro that each main file defines in its preamble
%% (English in notes-en.tex, Korean in notes-ko.tex).  The same file therefore
%% draws the same picture with the same number in both editions.  Never fork a
%% figure per language; add the new word to BOTH preambles instead.  Only the
%% \caption text lives in the language files.
%%
%% Black and grey only; curves are told apart by line style (solid, dashed,
%% dotted), never by colour.  Plotted curves are computed from the formula the
%% caption states, not sketched by hand.
%%
%% One macro per figure, \fig<Name>, in the order the figures first appear.
%% Large generated data (simulation snapshots, many-point curves) goes in its
%% own macro written by a script; say which script in a comment above it.

%% ===== (1) a diffusing profile spreads like a Gaussian =====================
%% c(x,t) = exp(-x^2/4Dt) / sqrt(4 pi D t), drawn for Dt = 0.25, 1, 4.
\newcommand{\figDiffusion}{%
\begin{tikzpicture}[x=0.9cm,y=4.2cm,>=Stealth]
  \draw[->] (-4.3,0) -- (4.5,0) node[right] {\lblPosition};
  \draw[->] (0,0) -- (0,1.32) node[left] {\lblConcentration};
  \foreach \x in {-4,-2,2,4} \draw (\x,0) -- (\x,-0.02) node[below] {\footnotesize $\x$};
  \draw[thick,samples=120,domain=-4.2:4.2]
    plot (\x,{exp(-\x*\x/1.0)/sqrt(3.14159265*1.0)});
  \draw[thick,dashed,samples=120,domain=-4.2:4.2]
    plot (\x,{exp(-\x*\x/4.0)/sqrt(3.14159265*4.0)});
  \draw[thick,dotted,samples=120,domain=-4.2:4.2]
    plot (\x,{exp(-\x*\x/16.0)/sqrt(3.14159265*16.0)});
  \node[anchor=west,font=\footnotesize] at (0.55,1.05) {$Dt=0.25$};
  \node[anchor=west,font=\footnotesize] at (1.25,0.36) {$Dt=1$};
  \node[anchor=west,font=\footnotesize] at (2.6,0.17) {$Dt=4$};
  \node[font=\footnotesize,align=center] at (-3.0,0.95) {\lblSpreads};
\end{tikzpicture}}

\endinput


\coursename{Sample Course}
\term{Fall 2027}
\lectitle{Lecture Notes, Lectures 1--2}
\lecdate{Updated 6 September 2027}
\notesby{Student}

\begin{document}
\makelecheader

% ===================================================================== front matter
\section*{Course information}
\begin{itemize}[leftmargin=1.2em,itemsep=1pt,topsep=3pt]
  \item Meets Mon/Wed. Midterm and final dates as announced in the syllabus.
  \item Grading as announced in the syllabus. Textbook: as announced.
\end{itemize}

\section*{The arc of the course}
\begin{itemize}[leftmargin=1.2em,itemsep=1pt,topsep=3pt]
  \item Part I --- Random walks and diffusion \hfill {\small[in progress]}
  \item Part II --- Steady states \hfill {\small[planned: announced in Lecture 2]}
\end{itemize}

\section*{How these notes are organised}
The notes are ordered by topic, not by date. A hairline carrying a lecture
number in the right margin marks the point where each lecture began, and the
grey line under it says what the notes for that lecture rest on. Material
added while writing the notes (a missing step, a check) is set small with a
$\blacktriangleright$ in the margin.

\section*{Lecture log}
{\small
\begin{longtable}{@{}l l l r@{}}
\toprule
lecture & date & sections & page\\
\midrule
\endhead
1 & Wed, 1 Sep 2027 & \S\ref{sec:walk} & \pageref{lec:2027-09-01}\\
2 & Mon, 6 Sep 2027 & \S\ref{sec:diffusion} & \pageref{lec:2027-09-06}\\
\bottomrule
\end{longtable}}

{\small\tableofcontents}

% ===================================================================== Part I
\partdivider{Part I --- Random walks and diffusion}

\lecturebreak[2027-09-01]{Lecture 1}
\sourcenote{handout 1 and the recording}
\section{The random walk}\label{sec:walk}

\begin{definition}\label{def:walk}
A \keyword{random walk} on a line is a sequence of positions
$x_n = \sum_{k=1}^{n} s_k$ whose steps $s_k=\pm a$ are independent and equally
likely.
\end{definition}

Each step has mean zero and variance $a^2$, and variances of independent steps
add, so
\begin{equation}\label{eq:var}
  \avg{x_n} = 0, \qquad \avg{x_n^2} = n a^2 .
\end{equation}

\scribe{The second result needs only that the steps are uncorrelated:
$\avg{s_j s_k}=0$ for $j\neq k$ kills every cross term of $\avg{x_n^2}$.}

% ---------------------------------------------------------------------------
\lecturebreak[2027-09-06]{Lecture 2}
\sourcenote{handout 2}
\section{The diffusion equation}\label{sec:diffusion}

With steps taken every $\tau$, write $t=n\tau$ and $D=a^2/2\tau$. Then
\eqref{eq:var} reads $\avg{x^2}=2Dt$, and in the continuum limit the density
of walkers obeys
\begin{equation}\label{eq:diffeq}
  \pd{c}{t} = D\,\frac{\partial^2 c}{\partial x^2}.
\end{equation}
Its solution for walkers that all start at the origin is
\begin{equation}\label{eq:gauss}
  c(x,t) = \frac{1}{\sqrt{4\pi D t}}\, e^{-x^2/4Dt},
\end{equation}
drawn in Figure~\ref{fig:diffusion}.

\begin{figure}[ht]
  \centering
  \figDiffusion
  \caption{The profile \eqref{eq:gauss} for $Dt=0.25$, $1$ and $4$ (solid,
  dashed, dotted). The area stays $1$ while the width grows as $\sqrt{2Dt}$.}
  \label{fig:diffusion}
\end{figure}

\begin{caution}[two conventions]\label{caut:dconv}
Some texts define $D=a^2/\tau$ and write $\avg{x^2}=Dt$. Check which one a
formula uses before combining it with \eqref{eq:diffeq}.
\end{caution}

% ===================================================================== Part II
\partdivider{Part II --- Steady states}

\lecturebreak[2027-09-08]{Planned}
\section{Steady states and boundary conditions}\label{sec:steady}

Written before the class from the posted outline; checked against the class
afterwards.

\phantomsection\label{lec:end}
\end{document}
